% =====================================================================
\chapter{Introducción}\label{cap:intro}
% =====================================================================

\section{Motivación y antecedentes}

En el curso de Óptica, la difracción suele presentarse primero a través de sus límites
clásicos —la difracción de campo lejano o de Fraunhofer, donde la propagación libre realiza una
transformada de Fourier espacial exacta, y la difracción de campo cercano o de Fresnel,
gobernada por el principio de Huygens--Fresnel— junto con las herramientas de la óptica de
Fourier y el filtrado espacial en el plano de frecuencias \parencite{goodman2005}. Sin embargo,
esos mismos fenómenos ondulatorios convierten al espacio libre y a los elementos difractivos en
procesadores analógicos masivamente paralelos: al propagarse entre planos modulados, un frente
de onda coherente ejecuta combinaciones lineales complejas a la velocidad de la luz y sin
consumo energético adicional al de la fuente de iluminación.

Mientras que el aprendizaje profundo convencional se ejecuta en procesadores electrónicos cuyo
costo energético y latencia crecen rápidamente con la dimensión de las matrices de pesos,
\textcite{lin2018} propusieron aprovechar la propagación difractiva en una \emph{red neuronal
difractiva profunda} (D\textsuperscript{2}NN): una pila de placas cuya variación local de
espesor impone retardos de fase independientes a la onda incidente. Entrenada numéricamente
en computadora mediante retropropagación y fabricada luego como objeto físico pasivo, la red
clasifica dígitos escritos a mano únicamente mediante interferencia y difracción.

Esa primera demostración operó en la banda de los terahertz ($\lambda\approx\SI{0.75}{mm}$),
donde neuronas de cientos de micrómetros se imprimen en 3D con facilidad. Estudios posteriores
profundizaron en la capacidad de representación de estas arquitecturas, su sensibilidad al
número de capas y las ventajas de integrarlas con etapas electrónicas de detección
\parencite{mengu2020}, confirmando que la profundidad óptica aporta grados de libertad que una
sola máscara no alcanza. No obstante, trasladar el esquema al espectro visible ofrece ventajas
decisivas: reduce el tamaño transversal del sistema en tres órdenes de magnitud y habilita
fuentes láser compactas y detectores CMOS convencionales. A cambio, exige fabricar neuronas de
escala micrométrica por fotolitografía y plantea una restricción física fundamental sobre el
espaciamiento entre planos: al disminuir $d_f$ o acortar $\lambda$, el cono de difracción de
cada píxel cambia de abertura y condiciona si las capas logran mezclarse entre sí.
\textcite{chen2021} abordan precisamente este régimen en el visible con un láser de He--Ne
($\lambda=\SI{632.8}{nm}$), proponen un criterio cuantitativo de conectividad entre capas y
demuestran experimentalmente un clasificador pasivo de cinco capas en sílice fundida.

\section{Planteamiento del problema}

El artículo de \textcite{chen2021} presenta su modelo físico (sus ecuaciones~1 a~5) y su
criterio de conectividad (ecuaciones~6 a~9) sin deducirlos, y valida el criterio con sólo
cuatro configuraciones numéricas. Tampoco describe cómo se discretiza la propagación ni
cómo se calcula el gradiente durante el entrenamiento. Este trabajo se propone llenar esos
huecos desde la teoría de la difracción escalar y comprobar los resultados mediante una
reproducción numérica independiente.

\section{Objetivos}

\paragraph{Objetivo general.} Deducir, justificar y reproducir numéricamente el modelo de
una D\textsuperscript{2}NN de solo fase en el visible propuesto por \textcite{chen2021}.

\paragraph{Objetivos específicos.}
\begin{enumerate}
  \item Deducir la función de transferencia del espacio libre (espectro angular) a partir de
        la ecuación de Helmholtz, y relacionarla con la integral de Rayleigh--Sommerfeld.
  \item Justificar el semiángulo de difracción $\varphi_{\max}=\arcsin[\lambda/(2d_f)]$ y
        deducir el criterio de conectividad entre capas.
  \item Obtener explícitamente el gradiente de la función de pérdida respecto a las fases.
  \item Analizar el muestreo numérico de la propagación y validar el propagador contra
        soluciones analíticas.
  \item Reproducir la Fig.~6 del artículo y explorar el criterio de conectividad con un
        barrido de la distancia entre capas.
  \item Estimar las tolerancias de fabricación del sistema.
\end{enumerate}

\section{Contribuciones}\label{sec:contribuciones}

Además de la reproducción, este trabajo aporta los siguientes elementos que no aparecen en el
artículo original:
\begin{enumerate}
  \item La deducción completa de las ecuaciones del modelo, desde la ecuación de Helmholtz
        hasta el criterio de conectividad (\cref{cap:difraccion,cap:modelo,cap:conectividad}).
  \item La justificación del factor~2 en $\varphi_{\max}=\arcsin[\lambda/(2d_f)]$ mediante el
        teorema de muestreo (\cref{sec:neurona}).
  \item El cálculo del gradiente con cálculo de Wirtinger y su interpretación como una onda
        de error que viaja en sentido inverso (\cref{sec:gradiente}).
  \item La demostración de que el criterio de conectividad coincide con la condición de
        muestreo del espectro angular de banda limitada (\cref{prop:identidad}), y de que su
        variable relevante es el número de Fresnel de una neurona (\cref{sec:fresnel-neurona}).
  \item Un barrido de la distancia entre capas que muestra que el criterio es unas diez veces
        más estricto de lo necesario (\cref{sec:barrido}).
  \item La identificación de dos inconsistencias en la Fig.~6 del artículo
        (\cref{sec:inconsistencias}) y un análisis de tolerancias de fabricación
        (\cref{sec:fabricacion}).
\end{enumerate}

\section{Organización del documento}

El \cref{cap:difraccion} desarrolla la teoría de la difracción escalar necesaria. El
\cref{cap:modelo} construye el modelo de la red y deduce su gradiente. El
\cref{cap:conectividad} analiza la conectividad entre capas. El \cref{cap:numerico} describe
la discretización y valida el propagador. El \cref{cap:resultados} presenta los resultados y
el \cref{cap:conclusiones}, las conclusiones. Los apéndices contienen la deducción de la
integral de Rayleigh--Sommerfeld y detalles de implementación.
